\documentclass[10pt,a4paper]{amsart}
\usepackage[margin=19mm]{geometry}
\usepackage{amsmath,amssymb,mathtools}
\usepackage[scr=rsfs,cal=euler]{mathalfa}
\usepackage{xfrac}
\usepackage[unicode,hidelinks,bookmarks=false]{hyperref}
\newcommand{\ie}{i.e.\ }
\newcommand{\cf}{cf.\ }
\newcommand{\resp}{respectively\ }
\newcommand{\Subsection}{\subsection}
\providecommand{\nobreakdash}{\nobreak\hbox{-}}
\setlength{\emergencystretch}{2em}
\multlinegap0pt
\usepackage{bm}
\usepackage{bbm}
\usepackage{dsfont}
\usepackage{xspace}
\usepackage{cancel}
\usepackage{mathtools}
\usepackage{amsthm}
\makeatletter
\@ifundefined{theorem}{\newtheorem{theorem}{Theorem}}{}
\makeatother
\title[Packing coloring of graphs with long paths]{Packing coloring of graphs with long paths}
\author{Hanna Furma\'nczyk, Didem G\"oz\"upek, Sibel \"Ozkan}
\date{Source: arXiv:2511.12761v1 (CC BY 4.0)}
\begin{document}
\maketitle

\noindent\emph{Source and licence.} Packing coloring of graphs with long paths. Version arXiv:2511.12761v1 (CC BY 4.0), distributed under the Creative Commons Attribution 4.0 International licence, as recorded in the arXiv licence metadata for this exact version and shown on the abstract page https://arxiv.org/abs/2511.12761v1. Full rights notice: licenses/arxiv-2511.12761-CC-BY-4.0.txt.\par\smallskip

\noindent\textit{Editorial scope.} Counted unit: the Theorem whose source label is \texttt{thm:caterpillar} in the article (source0.tex lines 1372--1375, source label \texttt{thm:caterpillar}), delivered together with the article's own complete original proof of it (source0.tex lines 1376--1510, one \texttt{proof} environment of 135 lines). No mathematics is written, repaired or re-derived: statement and proof are the article's own text line for line. Required context retained uncounted: none. Every definition, display and label of those lines is retained unchanged. Omitted from this package and not redistributed: 1--1371, 1511--1552. Nothing inside a retained excerpt is omitted or altered: every retained body is delivered exactly as the article's lines are written, so no omission receipt of any kind accompanies this package. Citations: the retained excerpts carry no citation command at all, so no bibliography entry of the article is transcribed and no reference is redistributed. Reference adaptation: none was needed and none was made; every reference inside the delivered unit resolves inside it, so no unresolved reference or citation is produced. Printed slips kept literal: no printed word, symbol, hypothesis or number of the counted statement or of its proof was altered. Layout and class: the article's own class is \texttt{article} with packages including graphicx, color, amsthm, amsmath, amsfonts, hyperref, todonotes, pgf, tikz, calc, rotating, tabu, authblk, slashbox; the wrapper is the lane's pinned \texttt{amsart} editorial wrapper at 10pt on A4, which restates the theorem-like environments the retained text needs and adds the article's own macro definitions that the retained text needs (none), with extra wrapper packages (bm, bbm, dsfont, xspace, cancel, mathtools). The delivered numbering is the unit's own. Assets: no retained span contains a figure environment, a graphics-inclusion command, a PDF-page inclusion, a table or a TikZ picture; no image file is copied into this package, the package declares no asset dependency and no image file is redistributed with it. Licence: version arXiv:2511.12761v1 of this article is distributed under the Creative Commons Attribution 4.0 International licence, as recorded in the arXiv licence metadata for this exact version and shown on the abstract page https://arxiv.org/abs/2511.12761v1. A full rights notice is delivered at licenses/arxiv-2511.12761-CC-BY-4.0.txt. Characters: every retained excerpt is pure ASCII, so the delivered engine, XeLaTeX with its default Latin Modern text font, reports no missing character.
\begin{theorem} \label{thm:caterpillar}
Let $G$ be a caterpillar $C(l;m_1,m_2,\ldots,m_l)$ with $l\geq 2$, $m_1,m_l \geq 1$, and $m_2,\ldots,m_{l-1}\geq 0$. Then, $\chi_p(G)=3$, if and only if $G \in \cup_{i=1}^{7}\mathcal{G}_i$. 

\end{theorem}

\begin{proof}
    The proof is based on the observation that the colors assigned to $v_1$ and $v_2$ largely constrain the coloring of the remaining vertices in the caterpillar $G$.

 ($\Rightarrow$) Let $G$ be a caterpillar of the form $C(l;m_1,m_2,\ldots,m_l)$ with $l\geq 2$, $m_1,m_l \geq 1$, and $m_2,\ldots,m_{l-1}\geq 0$ such that $\chi_p(G)=3$. Then, there are six possibilities to assign colors to $v_1$ and $v_2$. We will now analyze these six cases separately. In each case, we start by coloring the vertices $v_1$ and $v_2$. Then, we consider how this coloring can be extended to the remaining vertices of the backbone and to the leaves of $G$. In general, we assume that the vertices of the backbone that receive a color different from 1 may have neighbors of degree 1 corresponding to $m_i \geq 0$, for $2\leq i\leq l-1$, and $m_i \geq 1$, for $i=1$ or $i=l$. Vertices $v_i$ colored 1 may have adjacent leaves only in the case of $i=1$ or $i=l$, and in these cases exactly one leaf is allowed.
~~\\

\emph{Case 1:} Let $c(v_1)=1$ and  $c(v_2)=2$. 

First, suppose that $v_1$ has at least two leaves. Then, the leaves together with $v_2$ would imply that $\chi_p(G) \ge 4$, a contradiction. Moreover, by the definition of a caterpillar, $v_1$ is adjacent to at least one leaf; hence, $v_1$ is adjacent to exactly one leaf. Since $c(v_2)=2$, this leaf has to be colored with $3$. Then, $c(v_3)=1$ and $c(v_4)=3$ are the only possible assignments. It is easy to observe that the coloring pattern for the backbone becomes $[1\ 2\ 1\ 3]^{*}$. 
%\textcolor{blue}{It is easy to observe that the coloring pattern for the backbone becomes $1\ 2\ [1\ 3\ 1\ 2]^{*}$.}
Indeed, the only undetermined color in this sequence is a color of the neighbor of the vertex colored with 3; that is, the color of vertices $v_{4s+1}$, for $s\geq 1$. Here, the vertex can be colored with 2 or with 1. However, if color 2 is assigned to $v_{4s+1}$, then color 1 must be used for $v_{4s+2}$, which in turn necessitates the use of color 4 for the remaining vertices in the backbone. Therefore, this case is possible only at the end of the backbone where $l=4k+1$.
Hence, let us assume that we have colored $l - (l \bmod 4)$ vertices of the backbone of $G$ with the pattern $[1\ 2\ 1\ 3]^{*}$. Now, we need to consider how to color the rest of $l \bmod 4$ vertices of the backbone. 
%\textcolor{blue}{Hence, let us assume that we have colored $l - ((l-2) \bmod 4)$ vertices of the backbone of $G$ with the pattern $1\ 2\ [1\ 3\ 1\ 2]^{*}$. Now, we need to consider how to color the rest of $(l-2) \bmod 4$ vertices of the backbone.}

\underline{Let $l=4k$:} Here, the pattern $[1\ 2\ 1\ 3]^{*}$ colors all the vertices of the backbone. Hence, $G \in \mathcal{G}_{1}$.

%\textcolor{blue}{\underline{Let $l=4k$:} Here we have two uncolored vertices in the backbone and the only possible extension of the pattern to those vertices is to assign color 1 to $v_{l-1}$ and color 3 to $v_{l}$. In this case, we have $G\in \mathcal{G}_4$. }

\underline{Let $l=4k+1$:} The only vertex remaining uncolored in the backbone is $v_{4k+1}$, which may be assigned color 1 or 2. If color 1 is assigned to $v_{4k+1}$, then it may have only one adjacent leaf, which must be assigned color 2, resulting in a graph $G\in \mathcal{G}_2$; otherwise, the resulting graph is again in  $G\in \mathcal{G}_2$, where multiple adjacent leaves are permissible for $v_{4k+1}$.

%\textcolor{blue}{\underline{Let $l=4k+1$:} Since $c(v_{4k})=3$, it follows that $c(v_{4k+1})=1$, $c(v_{4k+2})=2$, and $c(v_{4k+3})=1$. Note that, the leaf adjacent to $v_{4k+3}$ must be assigned color 3. Indeed, the distance of this leaf to $v_{4k}$, which is also colored 3, is four; hence, the resulting coloring constitutes a feasible packing coloring. So, $G \in C(4k+3; 1, m_2, 0, m_4, 0, m_6, \ldots, 0, m_{4k+2},1)$, for any $k\geq 1$, and therefore belongs to the family $\mathcal{G}_5$.}

\underline{Let $l=4k+2$:} Now, we have two uncolored vertices in the backbone and the only possible extension of the pattern to those vertices is to assign color 1 to $v_{4k+1}$ and color 2 to $v_{4k+2}$.  
In this case, we have $G\in \mathcal{G}_4$.

\underline{Let $l=4k+3$:} If $k=0$, then we color the vertices of the backbone with colors $1,2,1$, consecutively. Note that, this is the only possible coloring since we have exactly one leaf adjacent to $v_1$ that is colored with 3. Moreover, $v_3$ has to be adjacent to exactly one leaf that is also assigned color 3. Therefore, $G \in C(3;1,m_2,1)$ which is in the family $\mathcal{G}_6$. 
Now, let $k\geq 1$. Since $c(v_{4k})=3$, it follows that $c(v_{4k+1})=1$, $c(v_{4k+2})=2$, and $c(v_{4k+3})=1$. Note that, the leaf adjacent to $v_{4k+3}$ must be assigned color 3. Indeed, the distance of this leaf to $v_{4k}$, which is also colored 3, is four; hence, the resulting coloring constitutes a feasible packing coloring. So, $G \in C(4k+3; 1, m_2, 0, m_4, 0, m_6, \ldots, 0, m_{4k+2},1)$, for any $k\geq 1$, and therefore belongs to the family $\mathcal{G}_5$.
    

~~\\
\emph{Case 2:} Let $c(v_1)=1$ and  $c(v_2)=3$.

Here, we follow a similar approach. The vertex $v_1$ may have only one leaf that must be assigned color 2. To extend the coloring to the backbone, we are required to use the coloring pattern $[1\ 3\ 1\  2]^*$. Indeed, the only undetermined color in this sequence is the color of the neighbor of the vertex colored with 3, namely, the color of vertices $v_{4s+3}$, for $s\geq 1$. Here, the vertex may be colored with 1 or 2. However, if color 2 is assigned to $v_{4s+3}$, then $v_{4s+4}$ must be colored with 1, which subsequently necessitates the use color 4 for the remaining vertices along the backbone. Therefore, such a situation can occur only at the end of the backbone, and moreover, only for graphs $G$ with $l=4k+3$.
\\Hence, let us assume that the initial $l - (l \bmod 4)$  vertices of the backbone of $G$ are colored using the pattern $[1\ 3\ 1\ 2]^*$. In the following, we consider how to extend the coloring for the remaining $l \bmod 4$ vertices of the backbone.

\underline{Let $l=4k$:} Here, the backbone is colored by the pattern $[1\ 3\ 1\ 2]^*$, yielding $C(4k;1,m_2,0,m_4,\ldots, 0, m_{4k}) \in \mathcal{G}_{1}$.

\underline{Let $l=4k+1$:} The only vertex  in the backbone remaining uncolored is $v_{4k+1}$, which must be assigned color 1, and it may have at most one adjacent leaf, which must be assigned color 3 (the distance to the nearest vertex colored with 3, that is to $v_{4k-2}$, is 4, yielding a feasible packing coloring). The family of caterpillars admitting this coloring can be described as follows: $C(4k+1; 1, m_2, 0, m_4, \ldots, m_{4k}, 1)$, for any $k\geq 1$, which is a member of $\mathcal{G}_2$.

\underline{Let $l=4k+2$:} Here we have two uncolored vertices in the backbone and the only possible extension of the pattern to those vertices is to assign color 1 to $v_{4k+1}$ and color 3 to $v_{4k+2}$. In this case, the family of caterpillars is described as $C(4k+2;1,m_2,0,m_4,0,m_6, \ldots,0,m_{4k+2}) \in  \mathcal{G}_4$, for any $k\geq 0$.

\underline{Let $l=4k+3$:} The three uncolored vertices may be assigned the colors $1, 3, 1$ or $1,3, 2$, respectively. It is easy to see that the latter coloring permits a larger number of leaves adjacent to $v_l$. In either extension, we have $G\in \mathcal{G}_5$, provided that $k\geq 1$.  If $k=0$, that is, if $l=3$, we likewise apply the coloring  $1, 3, 1$ or $1,3, 2$ to the vertices of the backbone; however, accounting for symmetry, this results in caterpillars belonging to the family $\mathcal{G}_6$.

~~\\
\emph{Case 3:} Let $c(v_1)=2$ and  $c(v_2)=1$. 

First, note that unlike in the previous two cases, vertex $v_1$ may have any number of leaves attached. To extend the coloring of $v_1$ and $v_2$ to the entire backbone we obtain the pattern $[2\ 1\ 3\ 1]^*$. Again, the only undetermined color in this sequence is the color of the neighbor of the vertex colored with 3, that is, the color of the vertices $v_{4s}$, for $s\geq 1$. Here, the vertex may be colored with 1 or 2. Nevertheless, if color 2 is assigned to $v_{4s}$, then $v_{4s+1}$ must be colored with 1, which subsequently necessitates the use of color 4 for the remaining vertices along the backbone. Thus, such a situation can occur only at the end of the backbone and, moreover, only for graphs $G$ with $l=4k$, where $k\geq 1$. Therefore, $v_{4s}$ must be assigned color 1, for any $ 1 \leq s\leq k -1$.
Hence, let us assume that we have colored $l - (l \bmod 4)$ initial vertices of the backbone of $G$ using the pattern $[2\ 1\ 3\ 1]^*$. In the following, we extend the coloring for the remaining $l \bmod 4$ vertices of the backbone.

\underline{Let $l=4k$:} Here, all vertices of the backbone are colored using either the pattern $[2\ 1\ 3\ 1]^*$ or $[2\ 1\ 3\ 1]^*[2\ 1\ 3\ 2]$. Both situations, taking symmetry into account, yield a caterpillar belonging to the family $\mathcal{G}_{1}$.

\underline{Let $l=4k+1$:} The vertex $v_{4k+1}$ must be assigned color 2 and may have an arbitrary number of adjacent leaves, which must be assigned color 1. Hence, $G\in \mathcal{G}_3$.

\underline{Let $l=4k+2$:} In this case, two vertices of the backbone remain uncolored, and the only possible extension of the pattern assigns color 2 to $v_{4k+1}$ and color 1 to $v_{4k+2}$. Note that $v_l$ may have only one adjacent leaf, which must be assigned color 3. The distance to the other vertices colored with 3 is, once again, consistent with a feasible packing coloring. Thus, $G \in C(4k+2; m_1, 0, m_3, 0, \ldots, m_{4k+1}, 1)$, for any $k\geq 0$, which belongs to the family symmetric to $\mathcal{G}_4$. 

\underline{Let $l=4k+3$:} The three uncolored vertices can only be assigned colors $2, 1, 3$, consecutively. In this case, $G\in \mathcal{G}_7$.

~~\\
\emph{Case 4:} Let $c(v_1)=2$ and $c(v_2)=3$. 

In this case, the extension of the backbone coloring differs substantially from that in the preceding three cases. Note that starting with colors $2$ and $3$, the sequence must continue with colors $1, 2, 1, 3, \ldots$ and the first vertex with undetermined color is $v_7$; in general, $v_{4s+7}$, for $s\geq 0$, that is, the vertex adjacent to the one colored with $3$. At this point, vertex $v_{4s+7}$ may be assigned color 1 or color 2. However, if it is colored with 2, the backbone must terminate at this vertex. In other words, such a coloring is feasible only at the end of the backbone, which occurs when $l=4k+3$, $k\geq 1$. Hence, let us assume that $l - ((l-2) \bmod 4)$ initial vertices of the backbone of $G$ are colored using the pattern $2\ 3\ [1\ 2\ 1\ 3]^*$, and $(l-2)\bmod 4$ vertices of the backbone remain uncolored. We now extend the coloring to these vertices.

\underline{Let $l=4k$:} 
%In addition, let $k\geq 1$. 
In this case, two vertices remain uncolored, and the extension is uniquely determined: we assign color 1 to $v_{l-1}$ and color 2 to $v_l$. It follows that $G \in C(4k; m_1, m_2,0, m_4,\ldots, 0, m_l)$, which is in $\mathcal{G}_1$.

\underline{Let $l=4k+1$:} 
%In addition, let $k\geq 1$. 
Here, three vertices remain uncolored, and the only possible extension assigns the colors $1,2,1$, consecutively. Note that $v_l$ can have only one leaf, which must be assigned color 3. Again, the distance to the nearest vertex colored 3, namely $v_{l-3}$, is 4; hence, the coloring is a feasible packing coloring.
In this case, $G\in C(4k+1; m_1,m_2,0,m_4,\ldots,0,m_{4k},1)$, which is in $\mathcal{G}_2$.

\underline{Let $l=4k+2$:} Here, all vertices of the backbone of $G$ are colored using the pattern $2\ 3\ [1\ 2\ 1\ 3]^*$, implying that $G\in \mathcal{G}_4$. 

\underline{Let $l=4k+3$:} In this case, one vertex remains uncolored, namely the vertex $v_l$, and the vertex $v_{l-1}$ is assigned color 3. If $k=0$, that is, the length of the backbone is 3, the vertex $v_l$ must be colored 1, yielding the coloring $2, 3, 1$. Note that $v_l$ may have only one adjacent leaf, which must be assigned color 2. Finally, we obtain a caterpillar from the family $\mathcal{G}_6$. If $k\geq 1$, two possible colors  arise for $v_l$: (i) when $c(v_l)=1$, the vertex $v_l$ may have only a single adjacent leaf, which must be colored 2; or (ii) when $c(v_l)=2$, the vertex $v_l$ may have one or more adjacent leaves. In both cases, $G \in \mathcal{G}_5$. 

~~\\
\emph{Case 5:} Let $c(v_1)=3$ and $c(v_2)=1$. 

Here, we will extend the coloring from $v_1$ and $v_2$ to the entire backbone; in particular, we will show that the coloring pattern $3\ [1\ 2\ 1\ 3]^*$ must be used. Indeed, the only vertex with an undetermined color in this pattern is the neighbor of the vertex colored with 3. Either color 1 or 2 may be used to color this vertex. Nevertheless, color 2 is permissible only at the end of the backbone, that is, when $l=4k+2$. Thus, we assume that the first $l-((l-1) \bmod 4)$ vertices of the backbone of $G$ are colored according to the pattern $3\ [1\ 2\ 1\ 3]^*$. At this stage, $(l-1)\bmod 4$ vertices at the end of the backbone remain uncolored.

\underline{Let $l=4k$:} 
%Let $k\geq 1$. 
Three vertices remain uncolored, and colors $1,2,1$ must be used consecutively. Observe that $v_l$ may have at most one adjacent leaf, and this leaf must be assigned color 3. The distance to the nearest vertex colored  3, namely $v_{l-3}$, is four; hence, the coloring is feasible. Thus, $G \in C(4k;m_1,0,m_3,0, \ldots, 0, m_{4k-1},1)$, which is in $\mathcal{G}_1$, taking symmetry into account.

\underline{Let $l=4k+1$:} 
%And let $k\geq 1$. 
All vertices of the backbone of $G$ are colored with the pattern $3\ [1\ 2\ 1\ 3]^*$, and $G \in \mathcal{G}_3$.

\underline{Let $l=4k+2$:} 
%And let $k\geq 0$. 
Here, one vertex remains uncolored, namely the vertex $v_l$, while vertex $v_{l-1}$ is assigned color 3. There are two possibilities: $(i)$ $v_l$ can be colored 1, in which case it is adjacent to exactly one leaf which must be assigned color 2; or $(ii)$ $v_l$ can be colored 2, in which case it may have any number of adjacent leaves, with $m_l \geq 1$. In both cases, $G_4\in C(4k+2; m_1,0,m_3,0,\ldots,m_{4k+1},m_{4k+2})$, which is in $\mathcal{G}_4$, taking symmetry into account.

\underline{Let $l=4k+3$:} When $k=0$, the colors 3, 1, 2 are assigned to $v_1$, $v_2$, and $v_3$, respectively. Otherwise, the two uncolored vertices $v_{l-1}$ and $v_l$ must be colored with 1 and 2, respectively. Thus, in both cases, $G \in \mathcal{G}_7$.

~~\\
\emph{Case 6:} Let $c(v_1)=3$ and $c(v_2)=2$. 

Observe that such a coloring cannot be extended. In fact, if $v_3$ existed, it would need to be assigned color 1, and its neighbor, whether a leaf or the next vertex on the backbone, would necessarily require color 4, contradiction. 
Hence, $G\in C(2; m_1,m_2)$, which is in $\mathcal{G}_4$.

\vspace{0.3cm}
 ($\Leftarrow$) It remains to show that every graph in $\bigcup_{i=1}^7 \mathcal{G}_i$ has packing chromatic number 3. For each family of graphs in $\bigcup_{i=1}^7 \mathcal{G}_i$, we provide a corresponding coloring pattern for the backbone. The leaves adjacent to the vertices colored 2 or 3 are assigned color 1. When a leaf is adjacent to a vertex colored 1 (that is, to $v_1$ or $v_l$), we explicitly indicate its color.

\begin{quote}
    
\begin{itemize}
        \item $\mathcal{G}_1=C(4k; m_1, m_2, 0,m_4,0, m_6, \ldots,0,m_{4k})$, for any $k\geq 1$. 
        
        
        Here, the coloring pattern for the backbone is $2\ 3\ [1\ 2\ 1\ 3]^{k-1}\ 1\ 2$.
        \item $\mathcal{G}_2=C(4k+1; 1, m_2,0,m_4,\ldots,0,m_{4k},m_{4k+1})$, for any $k\geq 1$.

        
        In this case, the backbone follows the coloring pattern $[1\ 2\ 1\ 3]^k \ 2$. Note that $v_1$ has only one adjacent leaf, which must be colored 3.
        \item $\mathcal{G}_3=C(4k+1; m_1, 0,m_3,\ldots, 0,m_{4k+1})$, for any $k\geq 1$.
    
        In this case, the coloring pattern of the backbone is $[3\ 1\ 2\ 1]^k\ 3$.
    \item $\mathcal{G}_4=C(4k+2; m_1, m_2,0,m_4,\ldots, 0,m_{4k+2})$, for any $k\geq 0$.
    
    The coloring pattern for the backbone is $2\ 3\ [1\ 2\ 1\ 3]^k$.
    \item $\mathcal{G}_5=C(4k+3; m_1, m_2, 0, m_4, \ldots, 0, m_{4k+2},m_{4k+3})$, for any $k\geq 1$.
    
    The coloring pattern for the backbone is $2\ 3\ [1\ 2\ 1\ 3]^k \ 2$.
    \item $\mathcal{G}_6=C(3; m_1,m_2,1)$.

    In this case, the coloring pattern $2\ 3\ 1$ is used for the backbone. Note that $v_3$ is adjacent to only one leaf, which must be assigned color 2.
    
    \item $\mathcal{G}_7=C(4k+3;m_1,0,m_3,\ldots, 0,m_{4k+3})$, for any $k\geq 0$.

    The coloring pattern for the backbone is $2\ 1\ 3\ [1\ 2\ 1\ 3]^k$.
    \end{itemize}

\end{quote}
 The proof is complete.
\end{proof}

\end{document}
